\chapter{A special transcendence theorem for the circle application}
\label{app:pi-transcendence}
This appendix closes the transcendence input used in the circle-squaring
application. Its algebraic ingredients are the characteristic-zero
splitting-field and fixed-field theorems proved earlier. Its analytic
ingredients are the elementary complex exponential, its addition and
derivative laws, Euler's identity $e^{i\pi}=-1$, the real fundamental
theorem of calculus applied to complex-valued functions componentwise,
and the elementary integral estimate on a real interval. Thus it does
not take a transcendence theorem as a premise. Real and complex
analysis underlying these elementary laws remain declared prior
foundations, just as the usual meaning of $\pi$ does.
The argument is the classical factorial-polynomial method of Lindemann;
for historical context, see Randy K. Schwartz, \emph{Pi is Transcendental:
Von Lindemann's Proof Made Accessible to Today's Undergraduates}
(2 February 2006), available through the university-hosted
\href{https://home.agh.edu.pl/rudol/Pi_TranscendentalLindemann.pdf}{expository paper}.
The proof below states all polynomial, integrality and estimate steps
used here, rather than taking that exposition as an unproved premise.

\section{An integer power-sum lemma}
\begin{lemma}\label{lem:integer-power-sums}
Let $Q(T)=\prod_{j=1}^m(T-\gamma_j)$ be a monic polynomial in
$\Z[T]$, with its complex roots listed with multiplicity. For every
integer $r\ge0$, $s_r=\sum_{j=1}^m\gamma_j^r$ is an integer.
Consequently $\sum_j A(\gamma_j)\in\Z$ for $A\in\Z[T]$.
\end{lemma}
\begin{proof}
The case $r=0$ gives $s_0=m$. Write
$U(t)=\prod_j(1-\gamma_jt)=1+c_1t+\cdots+c_mt^m$, where
the $c_j$ are integer coefficients of the reversed polynomial $Q$.
Differentiating the product gives the formal identity
\[
 -tU'(t)=U(t)\sum_{r\ge1}s_rt^r.
\]
No convergence is being assumed: for any coefficient of degree at
most $d$, replace each $(1-\gamma_jt)^{-1}$ by the finite geometric
sum through degree $d$ and compute modulo $t^{d+1}$; multiplication
and the product rule give exactly this identity in those degrees.
Equating the coefficient of $t^r$ yields
\[
 s_r+c_1s_{r-1}+\cdots+c_{\min(m,r-1)}s_{r-\min(m,r-1)}
       =-r c_r,
\]
where $c_r=0$ for $r>m$. Induction proves $s_r\in\Z$ for every
$r$. Expand $A$ as a finite integer linear combination of powers to
obtain the last assertion.
\end{proof}

\section{A relation among exponentials cannot have this form}
\begin{lemma}\label{lem:lindemann-special-relation}
Suppose $\beta_1,\ldots,\beta_m$ are nonzero complex numbers,
listed with multiplicities, such that
$P(X)=\prod_j(X-\beta_j)\in\Q[X]$. For every positive integer
$k$, the equality
\[
 k+\sum_{j=1}^m e^{\beta_j}=0
\]
is impossible.
\end{lemma}
\begin{proof}
Choose a positive integer $a$ divisible by the denominators of all
coefficients of $P$. Set $\gamma_j=a\beta_j$ and
\[
 Q(T)=a^mP(T/a)=\prod_j(T-\gamma_j)\in\Z[T].
\]
To check integrality explicitly, if the coefficient of $X^{m-r}$
in $P$ is $b_r$, the corresponding coefficient in $Q$ is $a^rb_r$,
an integer for $r\ge1$. Its constant coefficient $Q(0)$ is
nonzero, since none of the $\beta_j$ is zero.

Let $p$ be a prime not dividing the fixed nonzero integer
$kaQ(0)$. There are arbitrarily large such primes: otherwise the
finitely many primes in their union with the prime divisors of
$kaQ(0)$ would contain every prime, contrary to Euclid's proof
(a prime divisor of one plus their product is new). Put
\[
 H_p(T)=T^{p-1}Q(T)^p,\qquad N=(m+1)p-1,
 \qquad f_p(x)=\frac{H_p(ax)}{(p-1)!},\qquad
 F_p(x)=\sum_{r=0}^{N} f_p^{(r)}(x).
\]
The polynomial $H_p$ has integer coefficients and degree $N$.
For any integer-coefficient polynomial $H$, its $r$th derivative is
$r!$ times an integer-coefficient polynomial: this follows by
differentiating a monomial and using
$\frac{d^r}{dT^r}T^d=r!\binom dr T^{d-r}$.
For $r\ge p$, the integer $r!/(p-1)!$ is divisible by $p$.

At $x=0$, derivatives with $r<p-1$ vanish, because $f_p$ has a
factor $x^{p-1}$. The derivative at order $p-1$ equals
$a^{p-1}Q(0)^p$. All derivatives of order $r\ge p$ at zero are
integers divisible by $p$, by the preceding derivative formula.
It follows that
\[
 F_p(0)\in\Z,\qquad
 F_p(0)\equiv a^{p-1}Q(0)^p\pmod p.                 \tag{1}
\]
At each $\beta_j$, the multiplicity of its corresponding root
$\gamma_j$ in $H_p$ is at least $p$. Hence
$f_p^{(r)}(\beta_j)=0$ for $r<p$. For every $r\ge p$, the chain
rule and derivative formula give an integer-coefficient polynomial
$A_{p,r}$, independent of $j$, with
\[
 f_p^{(r)}(\beta_j)
   =\frac{a^rH_p^{(r)}(\gamma_j)}{(p-1)!}
   =p A_{p,r}(\gamma_j).
\]
Summing over $j$ and applying Lemma~\ref{lem:integer-power-sums}
shows
\[
 \sum_{j=1}^m F_p(\beta_j)\in p\Z.                \tag{2}
\]
Thus the quantity
\[
 J_p=kF_p(0)+\sum_jF_p(\beta_j)
\]
is an integer, and by (1)--(2) it is congruent modulo $p$ to
$ka^{p-1}Q(0)^p$, which is nonzero. In particular $|J_p|\ge1$.

We obtain the opposite conclusion analytically. Polynomial
differentiation telescopes to $F_p'-F_p=-f_p$. On the straight
segment $z=t\beta_j$, real differentiation gives
\[
 \frac{d}{dt}\bigl(e^{-t\beta_j}F_p(t\beta_j)\bigr)
       =-\beta_j e^{-t\beta_j}f_p(t\beta_j).
\]
Apply the real fundamental theorem of calculus to its real and
imaginary parts, then multiply by $e^{\beta_j}$, to find
\[
 e^{\beta_j}F_p(0)-F_p(\beta_j)
   =\beta_j\int_0^1 e^{(1-t)\beta_j}f_p(t\beta_j)\,dt.
\]
If the purported exponential relation held, summing this equation
would therefore give
\[
 -J_p=\sum_j\beta_j\int_0^1
           e^{(1-t)\beta_j}f_p(t\beta_j)\,dt.       \tag{3}
\]
Choose $R\ge1$ bounding all $|a t\beta_j|$ for $0\le t\le1$
and all $|\gamma_j|$. Since $H_p(T)=T^{p-1}\prod_j(T-\gamma_j)^p$,
on these finitely many segments
\[
 |f_p(t\beta_j)|
 \le \frac{R^{p-1}(2R)^{mp}}{(p-1)!}
 \le \frac{C^p}{(p-1)!},\qquad C=R(2R)^m.
\]
Furthermore $|e^{(1-t)\beta_j}|\le e^{|\beta_j|}$. The integral
estimate in (3) gives
\[
 |J_p|\le
 \left(\sum_j|\beta_j|e^{|\beta_j|}\right)
       \frac{C^p}{(p-1)!}\longrightarrow0.
\]
Indeed the ratio of consecutive terms $C^p/(p-1)!$ is $C/p$,
eventually at most $1/2$, so these terms tend to zero, also along
the arbitrarily large allowed primes. This contradicts $|J_p|\ge1$.
\end{proof}

\section{Transcendence of pi}
\begin{theorem}\label{thm:pi-transcendental}
The real number $\pi$ is transcendental over $\Q$.
\end{theorem}
\begin{proof}
Suppose otherwise and let $\alpha=i\pi$. The number $i$ is
algebraic because $i^2+1=0$, so closure of algebraic elements under
multiplication makes $\alpha$ algebraic. It is nonzero. Let
$\alpha_1=\alpha,\alpha_2,\ldots,\alpha_n$ be its conjugates
over $\Q$ in $\C$, and let $L$ be the splitting field of its
minimal polynomial in $\C$. Characteristic zero makes $L/\Q$
finite Galois; its fixed field is $\Q$. Its automorphisms permute
the conjugates.

Euler's identity gives $1+e^{\alpha_1}=0$, so
\[
 0=\prod_{j=1}^n(1+e^{\alpha_j})
   =\sum_{S\subseteq\{1,\ldots,n\}}
       \exp\!\left(\sum_{j\in S}\alpha_j\right).
\]
The second equality is the finite distributive expansion and the
addition law of the exponential. List the nonzero subset sums,
with their multiplicities, as $\beta_1,\ldots,\beta_m$. Let $k$
be the number of zero subset sums. The empty subset ensures
$k\ge1$, and the singleton $\{1\}$ ensures $m\ge1$. Thus
$k+\sum_je^{\beta_j}=0$.

Every automorphism of $L/\Q$ permutes subsets of the conjugates,
preserves whether a sum is zero, and hence permutes the multiset
of the nonzero $\beta_j$. Therefore every coefficient of
$P(X)=\prod_j(X-\beta_j)$ is fixed by the Galois group and lies
in $\Q$. All the hypotheses of
Lemma~\ref{lem:lindemann-special-relation} now hold, a contradiction.
\end{proof}
In particular $\sqrt\pi$ is transcendental: if it were algebraic,
its square would be algebraic. The constructibility theorem proves
that every coordinate constructible from rational unit data is
algebraic. A square of area $\pi$ times the unit square would have
side length $\sqrt\pi$; this gives the promised circle-squaring
obstruction with its transcendence premise proved here.
